\documentclass[a4paper,11pt,fleqn]{article}%fleqn pour aligner les equations à gauche
\usepackage{../../Styles/Cours}


\newcommand{\chnum}{Chapitre 12}
\newcommand{\chtitle}{Pression dans un liquide}
\newcommand{\dtype}{AE2}

\begin{document}
	\maketitle
La pression dans un liquide ne dépend pas du volume d’eau environnant mais de la profondeur considérée. On
peut se demander quelle relation relie la pression de l’eau et la profondeur.\\
Le montage expérimental est le suivant : la sonde d’un manomètre est plongée dans une éprouvette graduée.
On mesure la pression de l’eau à différentes profondeurs.

\begin{itemize}
	\item Remplir l’éprouvette de \SI{500}{mL} d’eau.
	\item Relever la pression atmosphérique indiquée par la sonde et la noter dans le tableau ci-dessous. Elle
	correspond à la pression à la surface, à l’altitude $z=0$ .
	\item Plonger l’embout de la sonde dans l’eau et, dès que la pression change, noter sa valeur et l’altitude
	correspondante dans le tableau. Attention, l’altitude est négative car on descend. L’axe $Oz$ est orienté
	vers le haut.
	\item Continuer à plonger la sonde. Relevez la pression et l’altitude correspondante dès que la pression
	change. Faire en sorte d’avoir environ une dizaine de mesures.
\end{itemize}

\begin{center}
	\begin{tabular}{|>{\centering}m{.15\linewidth}|>{\centering}m{.06\linewidth}|>{\centering}m{.06\linewidth}|>{\centering}m{.06\linewidth}|>{\centering}m{.06\linewidth}|>{\centering}m{.06\linewidth}|>{\centering}m{.06\linewidth}|>{\centering}m{.06\linewidth}|>{\centering}m{.06\linewidth}|>{\centering}m{.06\linewidth}|>{\centering\arraybackslash}m{.06\linewidth}|}
		\hline
		Altitude $z$ (\unit{\meter}) & & & & & & & & &&\\
		\hline
		Pression $P$ (\unit{\pascal}) & & & & & & & &&& \\
		\hline
	\end{tabular}
\end{center}

\begin{enumerate}
	\item Représenter la situation à l’aide d’un schéma muni d’un axe $Oz$.
	\item Recopier les valeurs mesurées con convertissant la pression en Pa.
	\begin{itemize}
		\item Ouvrir une page de tableau et recopier les valeurs de $z$ en \unit{\meter} et de $P$ en \unit{\pascal}.
		\item Tracer la courbe représentant $P$ en fonction de $z$.
	\end{itemize}
	\item La pression dans un liquide est-elle proportionnelle à l’altitude ? Justifier.
\end{enumerate}
On parle en général de profondeur plutôt que l’altitude quand on descend en dessous de $0$ sur l’axe $Oz$. Il ne
faut pas oublier le signe « - » devant la valeur de profondeur.
\begin{itemize}
	\item Afficher l’équation de la droite obtenue.
\end{itemize}
\begin{enumerate}[resume]
	\item Recopier l’équation de la droite et faire le lien entre cette équation et les grandeurs mesurées.
	\item A quoi correspond l’ordonnée à l’origine ?
\end{enumerate}
Sachant que la formule reliant la pression mesurée dans un liquide a pour formule : $P-P_0 =\rho\cdot g\cdot ( z_0-z )$ .

\begin{enumerate}[resume]
	\item Par association, déterminer une valeur approchée de $g$ pendant l’expérience.
\end{enumerate}

\end{document}